This work presents the design and development of a framework named "EDGAr", which integrates a Multimodal Retrieval-Augmented Generation (MRAG) engine designed to optimize document understanding within educational scenarios. The system is composed of hybrid indexing pipelines and extraction algorithms that perform fine-grained searches to verify and isolate structured textual and visual components before formulating responses. The data flow between modules occurs via vector representations, and the retrieved information feeds the client application "LuminIA Reader", which acts as the application layer. The project aims to prevent hallucinations caused by noise in long documents, ensuring dense contexts with low computational cost. The planned architectural tests seek to demonstrate the pipeline's efficiency, validating its functionality and viability for future implementations in a functional prototype. This architecture proposes itself as a relevant contribution to educational technology, enabling more precise interactions between the student and the educational material by mitigating the semantic hallucination problem.
% Separe as Keywords por ponto  

\keywords{Retrieval-Augmented Generation; Computer Vision; Educational Materials; Fine-grained Retrieval.}